\begin{abstract}
We prove the Gaussian propeller conjecture: for every finite measurable
partition of a Euclidean space, the sum of the squared lengths of its
Gaussian first moments is at most $9/(8\pi)$.
For dimension at least two and at least three cells, three planar
sectors of angle $2\pi/3$, extended by an orthogonal Euclidean factor,
attain the bound.
\end{abstract}

\section{Introduction}
\label{sec:introduction}
Let $\gamma_n$ denote the standard Gaussian probability measure on
$\mathbb R^n$, with density
$(2\pi)^{-n/2}\exp(-\|x\|^2/2)$.
For a measurable set $A\subseteq\mathbb R^n$, its Gaussian first moment is
$z(A)=\int_A x\,d\gamma_n(x)$. We call this unnormalized vector its
\emph{centroid}. Thus a centroid is not divided by the measure of its
cell. For positive integers $d,k$ and a measurable partition
$\mathcal A=(A_1,\ldots,A_k)$ of $\mathbb R^d$, define
\[
 \mathcal F(\mathcal A)=\sum_{i=1}^k\|z(A_i)\|^2.
\]
Partitions are understood up to Gaussian null sets, and empty cells are
allowed. The cell probabilities are unrestricted. All norms and inner
products below are Euclidean.

The propeller conjecture asks whether three planar sectors meeting at
angles $2\pi/3$, extended by an orthogonal Euclidean factor, maximize
this functional whenever $d\ge2$ and $k\ge3$. It arose in Khot and
Naor's study of approximate kernel clustering~\cite{KhotNaor}, where
Gaussian first moments govern approximation and Unique Games hardness.
Their formulation uses $k$ cells in $\mathbb R^{k-1}$, with unused
cells permitted. They proved that a maximizer can be chosen to be a
conical partition in the span of its centroids and determined the
three-cell value~\cite[Theorem~1.2 and Corollary~3.4]{KhotNaor}.

Heilman, Jagannath, and Naor proved the bound for every finite partition
of $\mathbb R^3$~\cite[Theorem~1.1]{HJN}. Their computer-assisted
argument treats the remaining four-cell conical configurations through
spherical geometry and a finite verification. We use that theorem as
an established input. The new estimates below show that an extremizer with five or more
active cells cannot have value greater than $9/(8\pi)$. Here an active
cell has positive Gaussian measure. This proves the conjecture in all
remaining dimensions.

\begin{theorem}\label{thm:main}
For all positive integers $d$ and $k$, every measurable partition
$(A_1,\ldots,A_k)$ of $\mathbb R^d$ satisfies
\[
 \sum_{i=1}^k\left\|\int_{A_i}x\,d\gamma_d(x)\right\|^2
 \le\frac9{8\pi}.
\]
Empty cells are allowed. The constant is sharp whenever $d\ge2$ and
$k\ge3$: three planar sectors of angle $2\pi/3$, multiplied by the
orthogonal complement of their plane, attain equality, with any further
cells empty.
\end{theorem}

\subsection{The attaining propeller}
\label{sec:example}
The attaining partition illustrates an important feature of the
problem: in every dimension at least two, only three cells are needed
once at least three labels are available.

Consider a sector $P$ in $\mathbb R^2$ of opening
$2\alpha$, symmetric about the first coordinate axis. Polar integration
gives
\[
 \int_P x\,d\gamma_2(x)
 =\frac1{2\pi}\left(\int_0^\infty r^2e^{-r^2/2}\,dr\right)
      \left(\int_{-\alpha}^{\alpha}(\cos\theta,\sin\theta)\,d\theta\right)
 =\frac{\sin\alpha}{\sqrt{2\pi}}(1,0).
\]
For $\alpha=\pi/3$, three such sectors partition the plane and their
squared centroid lengths sum to
$3\sin^2(\pi/3)/(2\pi)=9/(8\pi)$.
An orthogonal Gaussian factor contributes zero to every centroid, so
the construction has the same value in every dimension at least two.
The construction proves the attainment assertion in
Theorem~\ref{thm:main}.
Figure~\ref{fig:propeller-sectors} shows the planar partition and its
centroid directions.
